\originalchapter{作业一至五}
本章保留原作业的字母组与组内题号; $\ln$ 一律表示自然对数. 题面依据 Scott Sheffield 的 MIT 18.600 Fall 2019 作业重构; 每节附官方原件入口. A 组为 Ross 外书题, 按本书的范围限制仅登记书目定位, 不复制题面. 以下所有解答均为\textbf{本书补充推导; 官方未提供作业解答}. 原件中的提示用于理解题意, 不视为官方完整答案.

\originalsection{作业一: 排列、桥牌与彩票}
\href{https://ocw.mit.edu/courses/18-600-probability-and-random-variables-fall-2019/resources/mit18_600f19_pset1/}{官方作业一}. 原件在题目之前说明课程学习和合作要求; 本节集中呈现数学题.

\paragraph{A 组: 外书题登记.}
Sheldon Ross, \textit{A First Course in Probability}, 第8版, Pearson Prentice Hall, 2009, 第1章: A.1 为 Problem 10, 含(a)--(e); A.2 为 Theoretical Exercise 11. 原作业印有完整题面, 本书按范围限制排除其外书文字, 并非题面缺失.

\begin{exercise}\label{ps1-b1}
原题 B.1. 在集合 $\{1,\ldots,n\}$ 的 $n!$ 个置换 $\sigma$ 中, 有多少个只含一个循环? 即从1出发反复作用 $\sigma$, 依次经过所有元素后才返回1.
\end{exercise}
\begin{solution}
把循环固定写成 $(1,a_2,\ldots,a_n)$. 因为1的后继依次确定其余元素的位置, 每个单循环恰有一个这种表示. 剩下 $n-1$ 个元素可任意排列, 故共有 $(n-1)!$ 个. 对 $n=1$, 唯一恒等置换也算一个长度1的循环, 公式仍成立.
\end{solution}

\begin{exercise}\label{ps1-b2}
原题 B.2. 对 $1\le k\le n-1$, 有多少个置换恰有两个循环, 长度分别为 $k$ 和 $n-k$?
\end{exercise}
\begin{solution}
先选长度 $k$ 的循环使用的元素, 有 $\binom nk$ 种; 在这些元素上有 $(k-1)!$ 个循环, 在剩余元素上有 $(n-k-1)!$ 个. 若 $k\ne n-k$, 两个循环能由长度区分, 因而数量是
\[
\binom nk(k-1)!(n-k-1)!=\frac{n!}{k(n-k)}.
\]
若 $n=2k$, 选中的元素块与补集交换后得到同一个置换, 上式恰好把每个置换数了两次. 此时答案为 $n!/(2k^2)$. 这里除以2是因为两个循环没有先后次序, 并不是循环表示的旋转重复; 后一种重复已经由 $(k-1)!$ 处理.
\end{solution}

\begin{exercise}\label{ps1-b3}
原题 B.3. 有多少个置换是对合, 即对每个 $j$ 都满足 $\sigma(\sigma(j))=j$? 原提示建议先求恰含 $k$ 个长度2循环的个数, 再对 $k$ 求和.
\end{exercise}
\begin{solution}
一个循环若长度超过2, 作用两次不会使其中元素回到原处; 反之, 长度1和2的循环都满足要求. 因此对合由不动点和互不相交的配对组成. 恰有 $k$ 对时先选 $2k$ 个非不动点, 再将它们分成 $k$ 对. 把这些元素排成一列并相邻配对, 得到 $(2k)!$ 种排列, 但每对内的交换和各对间的排列不影响配对, 所以配对数为 $(2k)!/(2^k k!)$. 其余元素都固定. 总数为
\[
\sum_{k=0}^{\lfloor n/2\rfloor}\binom n{2k}\frac{(2k)!}{2^k k!}
=\sum_{k=0}^{\lfloor n/2\rfloor}\frac{n!}{(n-2k)!2^k k!}.
\]
$k=0$ 对应恒等置换, 不能遗漏.
\end{solution}

\begin{exercise}\label{ps1-c1}
原题 C.1. 标准52张牌含4种花色, 每种13张. 不计顺序的13张桥牌手牌共有 $\binom{52}{13}$ 种. 有多少手牌的13张牌全属同一花色?
\end{exercise}
\begin{solution}
选一种花色后, 必须把它的13张牌全取走, 没有进一步选择. 所以恰有4种.
\end{solution}

\begin{exercise}\label{ps1-c2}
原题 C.2. 在同样的桥牌手牌中, 有多少手牌恰好包含4种花色中的2种?
\end{exercise}
\begin{solution}
先选两种花色, 有 $\binom42=6$ 种. 从这两种花色的26张牌中选13张, 有 $\binom{26}{13}$ 种; 但全部来自第一种或第二种花色的两种手牌不符合要求. 因而答案为 $6[\binom{26}{13}-2]$. 每手符合要求的牌确定唯一的两花色集合, 所以不重计.
\end{solution}

\begin{exercise}\label{ps1-c3}
原题 C.3. 有多少13张桥牌手牌恰好包含3种花色?
\end{exercise}
\begin{solution}
先选3种花色, 有4种. 对选定的花色集合, 从39张牌中选13张. 令 $E_i$ 表示第 $i$ 种花色没有出现. 每个 $E_i$ 有 $\binom{26}{13}$ 种手牌, 每个 $E_i\cap E_j$ 恰有一种手牌, 而三种都不出现不可能组成13张. 容斥给出
\[
4\left[\binom{39}{13}-3\binom{26}{13}+3\right].
\]
加回的3对应三种单花色手牌, 它们此前各被减去了两次而不是一次.
\end{solution}

\begin{exercise}\label{ps1-c4}
原题 C.4. 有多少13张桥牌手牌包含全部4种花色, 即每种花色至少一张? 原提示建议把牌堆推广为 $k$ 种花色、每色13张, 用 $N_k(m)$ 表示恰含 $m$ 种花色的13张手牌数, 考察 $\sum_{m=1}^kN_k(m)=\binom{13k}{13}$ 及 $N_k(m)=\binom kmN_m(m)$.
\end{exercise}
\begin{solution}
从总数中排除缺至少一种花色的手牌. 缺指定一种、两种、三种花色的数量分别为 $\binom{39}{13}$、$\binom{26}{13}$、1. 四种都缺不可能. 因而所求为
\[
\binom{52}{13}-4\binom{39}{13}+6\binom{26}{13}-4.
\]
也可减去前3题的答案来验证. 提示中的第一式按实际出现的花色数分类; 第二式先选出现的 $m$ 色, 再在这些颜色内保证全出现. 这两式逐级求 $N_m(m)$, 等价于上面的容斥, 说明推广不是改变题目而是整理分类关系.
\end{solution}

\begin{exercise}\label{ps1-d}
原题 D. Powerball 的一注号码由 $\{1,\ldots,69\}$ 中不计顺序的5个白球号码及 $\{1,\ldots,26\}$ 中一个红球号码组成. 开奖均匀随机, 共有 $\binom{69}{5}26=292201338$ 种结果. 固定你的号码后, 求匹配恰1、恰2、恰3个白球的结果数, 此时不限制红球; 再求同时匹配红球且恰2个白球的结果数. 把各数量除以总数, 给概率的计算器数值, 并写一句对这些数值的评论. 可选用 Powerball 的 Wikipedia 页面核对.
\end{exercise}
\begin{solution}
若恰中 $r$ 个白球, 必须从自己的5个白球号码中取 $r$ 个, 再从其余64个号码中取 $5-r$ 个; 红球有26种自由选择. 所以数量为 $26\binom5r\binom{64}{5-r}$. 若红球必须中, 将因子26改为1. 各项是
\[
\begin{array}{c|r|c}
\text{匹配条件}&\text{结果数}&\text{概率}\\\hline
1\text{白, 红任意}&82598880&0.2826779664\\
2\text{白, 红任意}&10832640&0.0370725202\\
3\text{白, 红任意}&524160&0.0017938316\\
1\text{红}+2\text{白}&416640&0.0014258662
\end{array}
\]
例如两白球匹配约有3.71\%的概率, 但三白球已降到约0.179\%; 常遇到小程度的接近, 并不意味着大程度匹配也常见. 这些事件按恰好匹配数计, 不能把“至少”匹配的结果加入.
\end{solution}

\begin{exercise}\label{ps1-e1}
原题 E.1. 证明正态密度的归一化公式
$\int_{-\infty}^{\infty}(2\pi)^{-1/2}\ee^{-x^2/2}\dd x=1$.
原提示: 把无归一化因子的积分平方, 化为平面二重积分, 再用极坐标计算.
\end{exercise}
\begin{solution}
记 $I=\int_\R\ee^{-x^2/2}\dd x$. 此积分有限: 在 $|x|\ge1$ 时 $x^2\ge|x|$, 可用可积的 $\ee^{-|x|/2}$ 控制尾部. 被积函数非负, Tonelli 定理允许先写乘积为二重积分, 再在平面上作极坐标变换:
\[
I^2=\int_{\R^2}\ee^{-(x^2+y^2)/2}\dd x\dd y
=\int_0^{2\pi}\int_0^\infty\ee^{-r^2/2}r\dd r\dd\theta=2\pi.
\]
径向积分令 $u=r^2/2$ 后为 $\int_0^\infty\ee^{-u}\dd u=1$. 因 $I>0$, 取正平方根得 $I=\sqrt{2\pi}$, 除以此常数即为所证.
\end{solution}

\begin{exercise}\label{ps1-e2}
原题 E.2. 对 $\lambda>0$, 证明 $\sum_{k=0}^{\infty}\ee^{-\lambda}\lambda^k/k!=1$. 可用指数函数的 Taylor 展开.
\end{exercise}
\begin{solution}
指数函数的幂级数在每个实数处绝对收敛, 且 $\ee^\lambda=\sum_{k\ge0}\lambda^k/k!$. 乘以与求和指标无关的 $\ee^{-\lambda}$ 就得1. 非负性及总和为1也验证了这些项能作为概率质量函数; 当 $\lambda=0$ 时另定义为集中在0的分布.
\end{solution}

\begin{exercise}\label{ps1-e3}
原题 E.3. 对正整数 $n$ 及满足 $p+q=1$ 的数 $p,q$, 证明 $\sum_{k=0}^n\binom nkp^kq^{n-k}=1$. 原提示为展开 $(p+q)^n$.
\end{exercise}
\begin{solution}
$n$ 个因子相乘时, 取其中 $k$ 个因子的 $p$ 项与其余因子的 $q$ 项, 得到 $p^kq^{n-k}$; 选择这 $k$ 个位置有 $\binom nk$ 种. 按 $k$ 求和就是二项式展开, 故左边为 $(p+q)^n=1$. 恒等式只需 $p+q=1$; 解释为概率还须 $p,q\ge0$.
\end{solution}

\begin{exercise}\label{ps1-e4}
原题 E.4. 对整数 $n\ge0$, 证明 $\int_0^\infty x^n\ee^{-x}\dd x=n!$, 使用分部积分与归纳. 原件请读者记住 E 组四个公式并写下“Got it!”.
\end{exercise}
\begin{solution}
设 $I_n$ 为该积分. $I_0=1$. 对 $n\ge1$, 在 $[0,R]$ 上分部积分得到
\[
\int_0^R x^n\ee^{-x}\dd x=-R^n\ee^{-R}+n\int_0^R x^{n-1}\ee^{-x}\dd x.
\]
边界 $x=0$ 的项为0. 当 $R\to\infty$, 指数衰减快于任意固定次幂, 所以 $R^n\ee^{-R}\to0$; 若 $I_{n-1}$ 有限, 此式也证明 $I_n$ 有限并等于 $nI_{n-1}$. 从 $I_0=1$ 归纳, 得 $I_n=n!$. 四个归一化公式均已验证.
\end{solution}

\originalsection{作业二: 公理与假设}
\href{https://ocw.mit.edu/courses/18-600-probability-and-random-variables-fall-2019/resources/mit18_600f19_pset2/}{官方作业二}. 本节中的随机置换、骰子和随机选择按题意取均匀分布; 现实中的主观概率与风险中性概率不必相同.

\paragraph{A 组: 外书题登记.}
Ross 第8版第2章: A.1 为 Problem 25, A.2 为 Theoretical Exercise 10, A.3 为 Theoretical Exercise 20. 原件公开完整题面, 这里仅登记题号, 不复制外书文字.

\begin{exercise}\label{ps2-b1}
原题 B.1. 2019年9月12日, 民主党总统候选人提名合约在 PredictIt 与 Betfair 的报价如下, 单位为美分. YES 合约在相应候选人得到提名时支付100美分, 否则支付0; NO 支付相反. 例如 PredictIt 的 Warren YES 买价35, NO 买价66; 买入 NO 相当于以34美分卖出 YES, 买卖价差为1美分.
\begin{center}
\begin{tabular}{lrrrr}
\toprule
候选人 & PredictIt YES & PredictIt NO & Betfair YES & Betfair NO\\
\midrule
Elizabeth Warren&35&66&34.5&66.2\\
Joe Biden&27&74&21.7&79.2\\
Bernie Sanders&17&84&13.5&87.2\\
Andrew Yang&12&89&6.1&94.6\\
Kamala Harris&11&90&10.9&89.6\\
Pete Buttigieg&9&92&5.0&95.8\\
Hillary Clinton&5&96&3.6&96.7\\
Cory Booker&4&97&2.5&98.3\\
Beto O'Rourke&3&98&0.8&99.6\\
Julian Castro&2&99&1.0&99.1\\
Tulsi Gabbard&3&98&1.0&99.3\\
Tom Steyer&3&98&0.8&99.9\\
Amy Klobuchar&2&99&0.7&99.5\\
\bottomrule
\end{tabular}
\end{center}
暂忽略税、手续费、利息以及交易平台在结算前倒闭的风险. 原件举例: 买 PredictIt 的 Biden NO 与 Betfair 的 Biden YES, 成本95.7而到期必收100; 买全部13个 PredictIt NO, 成本1180, 因至多一人获得提名, 至少收1200. 在原件所述手续费制度下后者约收1188, 且只需保留最坏结算情况下足够的抵押资金. 除这些例子外, 从表中找出三种其他无风险获利方式.
\end{exercise}
\begin{solution}
同一候选人的 YES 和 NO 在同一提名事件上互补, 在题设的理想化模型中恰有一份支付100. 所以在两个平台分别买互补合约, 总成本若低于100, 差额就是确定利润. 三个例子分别为
\[
\begin{array}{c|c|c|c}
\text{候选人}&\text{PredictIt NO}&\text{Betfair YES}&\text{确定利润}\\\hline
\text{Sanders}&84&13.5&2.5\\
\text{Yang}&89&6.1&4.9\\
\text{Buttigieg}&92&5.0&3.0
\end{array}
\]
以 Yang 为例, 若其获提名, YES 支付100而 NO 支付0; 若未获提名, 支付互换, 都收100而只付95.1. 这三组都不同于原件的 Biden 示例, 不依赖对各候选人胜率的判断. 互补性要求两个平台对结算事件使用一致定义; 本题把这一点纳入给定合约模型, 并按题意不计交易摩擦.
\end{solution}

\begin{exercise}\label{ps2-b2}
原题 B.2. 有效市场假说认为套利机会应被交易消除. 对表中看似存在的机会给出一个简短的推测解释; 原件只检查是否认真思考, 不按某个唯一解释评分. 原评注说 PredictIt 在当时限制每位交易者对每个候选人投入850美元, 其法律安排与 Betfair 不同. 例如2美分的 Klobuchar YES 可买42500份, 而99美分的 NO 在限额内只能买858份. 报价须以整数美分变化; 以较低价格挂单还涉及排队和成交失败的风险.
\end{exercise}
\begin{solution}
一个可能的解释是两市场受不同准入条件、仓位限制和执行成本影响, 所以套利资金不能无限进入. 即使每份合约有几美分价差, 投资限额也可能使一名交易者只能赚很少的钱, 不足以抵偿实际手续费、资金占用和操作成本. 题设为了计算删去了这些成本, 但市场报价仍可能由包含这些成本的现实交易者形成. 这是一种推测而非利用表格证明的事实, 也不是对所有时期报价的结论.
\end{solution}

\begin{exercise}\label{ps2-c}
原题 C. 有 $M\ge1$ 个求职者和 $N\ge1$ 家招聘公司. 每个求职者独立、均匀地对一家公司的职位产生认真兴趣; 每家公司也独立、均匀地对一个求职者产生认真兴趣, 所有选择相互独立. 求至少出现一对互相感兴趣的公司与求职者的概率, 可写为不再化简的和. 原提示为对事件 $E_j=$“第 $j$ 名求职者的兴趣得到回应”使用容斥. 再计算 $\sum_{j=1}^M\Prob(E_j)$, 即有回应的求职者人数的期望, 并评论是否出乎意料.
\end{exercise}
\begin{solution}
单个求职者先选任意公司, 该公司回应他的概率恒为 $1/M$, 所以 $\Prob(E_j)=1/M$. 但多个求职者的成功事件不独立: 同一公司只能回应一个人. 这正是使用容斥而不能直接相乘的原因.

固定 $k$ 个求职者. 若他们全得到回应, 他们选的公司必须互异. 有序选择互异公司的方法数为 $(N)_k=N(N-1)\cdots(N-k+1)$, 每个指定选择的概率是 $N^{-k}$; 然后这 $k$ 家公司分别选对人, 概率是 $M^{-k}$. 当 $k>N$ 时交事件为空. 因此
\[
\Prob\!\left(\bigcup_{j=1}^M E_j\right)
=\sum_{k=1}^{\min(M,N)}(-1)^{k+1}\binom Mk\frac{(N)_k}{N^kM^k}.
\]
设 $S=\sum_j\ind_{E_j}$. 指示变量线性相加给出
$\E S=\sum_j\Prob(E_j)=M/M=1$.
公司数即使增大, 每个求职者已选公司的回应概率也仍是 $1/M$, 所以此期望不依赖 $N$. 这有些出人意料, 但不表示一定有一对: $S$ 可以为0, 也可以大于1. 边界上, $M=1$ 或 $N=1$ 时恰有一对, 容斥式也给1.
\end{solution}

\begin{exercise}\label{ps2-d}
原题 D. Alice 与 Bob 的网球局达到平分(deuce). 此后逐分继续, 直到一人的总得分比另一人多2分时该人胜出. Alice 每分独立以概率 $p$ 获胜, Bob 以 $q=1-p$ 获胜. 求 Alice 赢下整局的概率. 原提示建议把后面两分分成 Alice 连胜、Bob 连胜与各胜一分三类. 判断: 若 Alice 赢一分的可能性为 Bob 的 $k$ 倍, 她在平分时赢整局的可能性是否为 Bob 的 $k^2$ 倍?
\end{exercise}
\begin{solution}
把两分视为一轮. Alice 赢下一轮并结束的概率是 $p^2$, Bob 是 $q^2$, 返回平分的概率是 $2pq$. 返回后, 未来独立且状态相同. 因为 $2pq\le1/2$, 一直返回的概率为 $\lim_{r\to\infty}(2pq)^r=0$, 所以比赛几乎必然结束. 按结束前返回次数分类,
\[
\Prob(\text{Alice胜})=\sum_{r=0}^\infty(2pq)^rp^2
=\frac{p^2}{1-2pq}=\frac{p^2}{p^2+q^2}.
\]
同理 Bob 的概率为 $q^2/(p^2+q^2)$. 对 $0<p<1$, 两者比为 $(p/q)^2=k^2$, 故陈述成立. $p=0$ 或1时一方必胜, 无需使用有限比值 $k$.
\end{solution}

\begin{exercise}\label{ps2-e}
原题 E. xkcd 的“随机”按钮从恰好2200幅已编号漫画中独立均匀选一幅. 点击 $m$ 次后, 至少有一幅被看到两次的概率是多少? 给出 $m=36,56,78$ 的粗略数值. 可借计算器计算 $\prod_{k=0}^{35}(1-k/2200)$. 判断: 首次看到重复漫画所需的点击数是一个随机量, 其中位数约56, 且约有一半概率落在36与78之间.
\end{exercise}
\begin{solution}
直接统计所有重复方式容易重叠, 所以先算一次都不重复的概率. 第 $k+1$ 次点击时, 若此前 $k$ 幅互异, 避开它们的条件概率是 $1-k/2200$. 因而对 $0\le m\le2200$,
\[
Q_m=\Prob(\text{前}m\text{次互异})=\prod_{k=0}^{m-1}\left(1-\frac{k}{2200}\right),
\qquad \Prob(\text{出现重复})=1-Q_m.
\]
$m>2200$ 时重复概率为1. 数值依次为 $0.250176,\ 0.506366,\ 0.748753$, 即约25\%、51\%、75\%.

设 $T$ 为首次重复的点击数, 则 $\Prob(T>m)=Q_m$. $1-Q_{55}<1/2\le1-Q_{56}$, 所以按最小满足分布函数达到一半的整数定义, 中位数为56. 又
\[
\Prob(36\le T\le78)=\Prob(T>35)-\Prob(T>78)=Q_{35}-Q_{78}\approx0.510699.
\]
端点若改为严格不等式, 数值稍变, “约一半”的结论不变. 所以原陈述的两个部分均可接受.
\end{solution}

\begin{exercise}\label{ps2-f}
原题 F. 一个公布新生儿性别的聚会上, 15个盖好的杯子装有彩色珠子. 若是男孩, 有8蓝7粉; 若是女孩, 有7蓝8粉. 以均匀随机顺序逐杯翻开, 直到看见8杯同色, 从而确定性别. 令 $N$ 为翻开的杯数. 求 $\Prob(N=k)$, $k=8,9,\ldots,15$, 并判断必须等到最后一杯的概率是否超过 $1/2$.
\end{exercise}
\begin{solution}
不论颜色对应哪种性别, 停止都发生在8个多数色杯子全部出现之时. 因而 $N$ 就是多数色杯子所占8个位置的最大值. 这8个位置在 $\binom{15}{8}$ 个集合中均匀分布. 最大值为 $k$ 时, 位置 $k$ 必须被占, 其余7个位置可从 $1,\ldots,k-1$ 中任选, 所以
\[
\Prob(N=k)=\frac{\binom{k-1}{7}}{\binom{15}{8}},\quad 8\le k\le15.
\]
对 $k$ 求和得到1, 因为每个8元素位置集合恰有一个最大值. 特别地 $\Prob(N=15)=\binom{14}{7}/\binom{15}{8}=8/15>1/2$. 也可直接看最后一杯是多数色的概率为8/15, 此时且仅此时需等到最后.
\end{solution}

\begin{exercise}\label{ps2-g}
原题 G. 均匀洗乱一副含26红26黑的52张牌. 主持人从顶端逐张翻牌供你观察, 你须在全部52张翻完之前说“现在”; 下一张若红则赢1美元, 若黑则无奖励. 你希望使获奖概率最大, 应采用怎样的停止策略? 原件提示考虑说“现在”后改翻最底下一张的变体, 比较两游戏在同一策略下的获奖概率.
\end{exercise}
\begin{solution}
策略可以依赖已翻开的牌, 也可使用与牌堆独立的私人随机数, 但不能预知未翻的牌. 停止时设还有 $r\ge1$ 张未翻, 其中 $a$ 张红. 条件于已见信息与停止决定, 未见牌的相对次序仍均匀; 其顶张与底张为红的条件概率都为 $a/r$. 因此按完全相同的观察策略停止, 在原游戏与底牌变体中获胜的总概率相同, 都是停止时 $a/r$ 的期望.

在底牌变体中, 策略只影响何时翻底牌, 不影响作为奖励依据的那张固定底牌. 又因必须保留至少一张, 底牌在决定之前一直未被翻开. 所以胜负事件就是“原牌堆底牌为红”, 概率恒为26/52. 从而原游戏中任何合法策略的获胜概率都为 $1/2$; 立即说“现在”、等待某个红牌比例或直到只剩一张, 都同样最优. 条件红牌比例有时会升高, 但等待也有降低它的风险, 两者在总体上正好抵消. 这给出了有限牌堆的完整论证, 不必预先调用鞅的可选停止定理.
\end{solution}

\originalsection{作业三: 条件概率}
\href{https://ocw.mit.edu/courses/18-600-probability-and-random-variables-fall-2019/resources/mit18_600f19_pset3/}{官方作业三}. 条件概率在条件事件概率为正时定义为 $\Prob(A\mid B)=\Prob(A\cap B)/\Prob(B)$. Bayes 公式只是在同一联合概率上交换条件方向; 数据怎样被选择, 也必须属于概率模型.

\paragraph{A 组: 外书题登记.}
Ross 第8版第3章, A.1 为 Problem 47. 原件含完整题面及两个无编号输出, 本书只保留这一定位, 不复制外书文字.

\begin{exercise}\label{ps3-b}
原题 B. 某诊所使用快速流感检测. 灵敏度为0.5, 即有流感者一半检测阳性; 特异度为0.9, 即无流感者九成检测阴性. 医生根据季节与症状, 对当前病人有流感的先验概率取 $\Prob(F)=0.7$, 并认为上述检测参数适用于该病人. 令 $T$ 为检测阳性, 所以 $\Prob(T\mid F)=0.5$, $\Prob(T\mid F^c)=0.1$. 求检测阴性后的 $\Prob(F\mid T^c)$ 以及阳性后的 $\Prob(F\mid T)$, 给近似百分数. 原件评注建议另查实际检测灵敏度、特异度, 并提醒药物与疫苗效果随研究、年份及地区变化; 这里仅计算给定模型.
\end{exercise}
\begin{solution}
先把检测结果的总体概率按是否有流感分解. 阴性概率为 $0.7(1-0.5)+0.3(1-0.1)=0.62$, 阴性且有流感的概率为0.35. 因此
\[
\Prob(F\mid T^c)=\frac{0.35}{0.62}=\frac{35}{62}\approx56.45\%.
\]
阳性的总概率为 $0.7\cdot0.5+0.3\cdot0.1=0.38$, 所以
$\Prob(F\mid T)=0.35/0.38=35/38\approx92.11\%$.
阴性会降低原先70\%的判断, 但不会降到很低: 检测漏掉一半患者, 而病人原先又很可能有病. 这不是把检测误差当成后验概率; 后验还取决于先验和两种人群的比例.
\end{solution}

\begin{exercise}\label{ps3-c1}
原题 C.1. 独立地无限次抛公平硬币, $X_i\in\{H,T\}$ 为第 $i$ 次结果. 已知前20次恰有9次正面, 求前10次恰有7次正面的条件概率.
\end{exercise}
\begin{solution}
给定总正面数为9后, 正面位置在20个位置的所有9元素子集中均匀. 前10个位置选7个正面, 后10个位置需选2个正面, 故
\[
\Prob\left(\sum_{i=1}^{10}\ind_{\{X_i=H\}}=7\ \middle|\ \sum_{i=1}^{20}\ind_{\{X_i=H\}}=9\right)
=\frac{\binom{10}{7}\binom{10}{2}}{\binom{20}{9}}.
\]
分母不是 $2^{20}$, 因为条件已删去了所有总正面数不等于9的序列.
\end{solution}

\begin{exercise}\label{ps3-c2}
原题 C.2. 求存在一个无限等差数列, 其所有位置的硬币结果全为正面的概率. 即存在正整数 $a,b$, 使 $X_i=H$ 对每个 $i\in\{a,a+b,a+2b,\ldots\}$ 成立. 原提示为可数可加性.
\end{exercise}
\begin{solution}
固定一对 $a,b$. 前 $r$ 个指定位置全为正面的概率是 $2^{-r}$. 这些事件随 $r$ 递减, 所以由概率对递减事件列的连续性, 无限全正面的事件 $E_{a,b}$ 的概率为 $\lim_r2^{-r}=0$. 正整数对 $(a,b)$ 只有可数多个, 故
\[
\Prob\left(\bigcup_{a,b\ge1}E_{a,b}\right)\le\sum_{a,b\ge1}\Prob(E_{a,b})=0.
\]
必须先指出候选数列可数; 不能对任意不可数个零概率事件照搬这一推理.
\end{solution}

\begin{exercise}\label{ps3-c3}
原题 C.3. 求固定模式 HHTTHTTHT 在 $X_1,X_2,\ldots$ 中至少出现一次的概率.
\end{exercise}
\begin{solution}
把硬币序列切成互不重叠的长度9区块. 每块恰为指定模式的概率是 $2^{-9}$, 各块事件独立. 前 $r$ 块都不出现该模式的概率为 $(1-2^{-9})^r\to0$. 若模式在整个序列中从不出现, 它尤其不会在这些区块起点出现, 所以其概率也为0. 至少出现一次的概率是1. 我们没有假定相邻起点的模式事件独立; 只使用了特意选择的互不重叠区块.
\end{solution}

\begin{exercise}\label{ps3-c4}
原题 C.4. 求每一种有限长度硬币模式都在序列中出现无穷多次的概率.
\end{exercise}
\begin{solution}
先固定长度 $\ell$ 的某个模式 $w$. 在不重叠的长度 $\ell$ 区块中, 对任何固定区块序号 $r$, 其后永远不再出现 $w$ 的概率为
$\lim_{m\to\infty}(1-2^{-\ell})^m=0$.
仅出现有限次意味着存在某个 $r$, 此后再无出现; 对 $r$ 作可数并, 此事件仍为零概率. 因而固定 $w$ 几乎必然出现无穷多次.

所有有限模式的集合是 $\bigcup_{\ell\ge1}\{H,T\}^{\ell}$, 为可数集合. 对每种模式“不出现无穷多次”的零概率例外再作可数并, 概率仍为0. 故所有模式同时出现无穷多次的概率为1. 第二次可数并确保结论是对所有模式同时成立, 而不只对事先指定的一种成立.
\end{solution}

\begin{exercise}\label{ps3-d}
原题 D. 审讯星球有730名嫌疑人, 恰有一人有罪. 有罪者对任一问题以0.9概率给出听来可疑的回答, 以0.1概率给出正常回答; 无罪者的概率相反. 条件于某人的有罪或无罪身份, 不同问题的回答相互独立. 均匀随机选一人询问9个问题, 前3个回答正常, 后6个连续可疑. 审讯者认为无罪者连续6次可疑的概率仅百万分之一, 故认定有罪. 给定实际9个回答, 求此人有罪的精确条件概率.
\end{exercise}
\begin{solution}
设 $G$ 为被选人有罪, $D$ 为观察到的有序回答序列. 先验为 $\Prob(G)=1/730$, $\Prob(G^c)=729/730$. 由题设条件独立,
\[
\Prob(D\mid G)=0.1^3 0.9^6,\qquad
\Prob(D\mid G^c)=0.9^3 0.1^6.
\]
二者的比为 $(0.9/0.1)^3=9^3=729$. Bayes 公式得到
\[
\Prob(G\mid D)=\frac{0.1^3 0.9^6}{0.1^3 0.9^6+729\,0.9^3 0.1^6}=\frac12.
\]
六个可疑回答确实支持有罪, 但前三个正常回答支持无罪, 且先验中无罪的可能性为有罪的729倍. 观察信息的似然比恰好抵消这一先验比. 忽略前3答或把 $\Prob(D\mid G^c)$ 直接当成 $\Prob(G^c\mid D)$, 都会改变所问概率.
\end{solution}

\begin{exercise}\label{ps3-e}
原题 E. 若 $\Prob(A\mid X_1),\ldots,\Prob(A\mid X_k)$ 均相等, 证明 $\Prob(X_i\mid A)$ 与 $\Prob(X_i)$ 成正比, 即比值不依赖 $i$. 不要求各 $X_i$ 互斥. 所有出现的条件概率都假定有定义.
\end{exercise}
\begin{solution}
记共同值为 $c$. 因 $\Prob(X_i)>0$ 且 $\Prob(A)>0$, Bayes 公式给出
\[
\frac{\Prob(X_i\mid A)}{\Prob(X_i)}
=\frac{\Prob(A\mid X_i)}{\Prob(A)}=\frac c{\Prob(A)}.
\]
这不依赖 $i$. 每个等式只使用 $A$ 与一个 $X_i$ 的交事件, 未对 $X_i$ 作分割或求和, 所以无需互斥. 若这些假设同样解释观测 $A$, 相对可信度不会因 $A$ 而改变, 原先的先验比例保留.
\end{solution}

\begin{exercise}\label{ps3-f}
原题 F. 保守科学星球上, 专家提出的假说以0.9概率为真, 记为 $\Prob(H)=0.9$. 公布前做一次实验, 预先定义阳性事件 $T$, 且 $\Prob(T\mid H)=0.95$, $\Prob(T\mid H^c)=0.05$. 只有阳性结果才发表.
\begin{enumerate}[label=(\alph*)]
\item 求 $\Prob(H\mid T)$, 即公开结果的可靠率.
\item 若第一次阳性, 再作一次同类实验, 其阳性事件记为 $\widetilde T$. 设 $\Prob(\widetilde T\mid H T)=0.95$, $\Prob(\widetilde T\mid H^c T)=0.05$. 求复现率 $\Prob(\widetilde T\mid T)$.
\item 求 $\Prob(H\mid T\widetilde T)$, 即复现后结果的可靠率.
\end{enumerate}
推测科学星球则先提出若为真会很有趣的假说, 但 $\Prob(H)=0.05$, 实验参数为 $\Prob(T\mid H)=0.8$, $\Prob(T\mid H^c)=0.05$. 对该星球:
\begin{enumerate}[label=(\alph*),start=4]
\item 求 $\Prob(H\mid T)$.
\item 求复现率 $\Prob(\widetilde T\mid T)$, 其中 $\Prob(\widetilde T\mid HT)=0.8$, $\Prob(\widetilde T\mid H^cT)=0.05$.
\end{enumerate}
原件还提出课外思考: 两种制度有何利弊? 如果使用者理解其含义, 较低可靠率与复现率是否总是坏事? 医学研究昂贵且数据稀缺时, 怎样改进决策? 这些思考不要求唯一数值答案.
\end{exercise}
\begin{solution}
先对任一先验 $h$, 真假阳性率 $s,f$ 写出一次阳性的概率
$d=hs+(1-h)f$. 发表流程把研究集合限制在事件 $T$ 上, 所以(a)(d)问的是 $hs/d$, 而不是无条件先验 $h$.

(a) 保守星球的 $d=0.9\cdot0.95+0.1\cdot0.05=0.86$, 从而
$\Prob(H\mid T)=0.855/0.86=171/172\approx99.42\%$.

(b) 条件于第一次阳性, 再按 $H,H^c$ 分类, 得
\[
\Prob(\widetilde T\mid T)=s\Prob(H\mid T)+f\Prob(H^c\mid T)
=\frac{hs^2+(1-h)f^2}{d}.
\]
因此保守星球的复现率为
$(0.9\cdot0.95^2+0.1\cdot0.05^2)/0.86
=0.8125/0.86=325/344\approx94.48\%$.
这里没有声称两次实验无条件独立. 题目规定第二次实验在已知假说真伪及第一次阳性后保持同样阳性率, 这足以在链式概率中相乘. 若两次实验还受共同未建模因素影响, 题给的条件概率就可能不成立.

(c) 两次阳性且 $H$ 为真的概率为 $hs^2=0.81225$, 两次阳性的总概率为0.8125. 故
$\Prob(H\mid T\widetilde T)=0.81225/0.8125=3249/3250\approx99.969\%$.

(d) 推测星球的 $d=0.05\cdot0.8+0.95\cdot0.05=0.0875$. 所以
$\Prob(H\mid T)=0.04/0.0875=16/35\approx45.71\%$.

(e) 使用同一个条件分解,
\[
\Prob(\widetilde T\mid T)
=\frac{0.05\cdot0.8^2+0.95\cdot0.05^2}{0.0875}
=\frac{0.034375}{0.0875}=\frac{11}{28}\approx39.29\%.
\]
较低先验使假阳性在已发表结果中占很大比例; 5\%的假阳性率并不意味着95\%的已发表结果为真. 对课外思考, 一种合理观点是探索阶段可容许较低可靠率来发现意外方向, 但应用于重大决策之前须以更强证据提高可靠率; 这并不由上述五个数值唯一决定, 因为还要比较错误和遗漏的成本. 原评注所列实际复现研究只作进一步阅读, 这些理想化数值不冒充那些研究的估计.
\end{solution}

\begin{exercise}\label{ps3-g}
原题 G. 第一家工厂生产10000批棒球卡, 每批100张, 编号1--100; 第二家生产1000批, 每批1000张, 编号1--1000; 第三家生产100批, 每批10000张, 编号1--10000. 三家总产量各100万张. 随机见到一张卡, 先验认为三家同样可能, 并在每家来源下对其卡均匀抽样. 观察到编号74.
\begin{enumerate}[label=(\alph*)]
\item 求来自第一、第二、第三家工厂的三个条件概率.
\item 类比地, 三个宇宙各含 $10^{50}$ 个智慧个体, 分别按每文明 $10^{12},10^{15},10^{18}$ 个个体分组. 从总共 $3\cdot10^{50}$ 个个体中均匀选一人, 得知在其出生前, 该文明恰有141452234521个个体出生. 求此人来自第一个宇宙的条件概率.
\end{enumerate}
原件以此引出末日论证, 但没有要求把这三宇宙模型直接当成人类文明的实际预测.
\end{exercise}
\begin{solution}
(a) 设 $F_i$ 为第 $i$ 家来源. 每家先验为1/3, 看到74的似然分别为 $1/100,1/1000,1/10000$. 批数已使总量相同, 但不改变各批内号码的相对频率. Bayes 公式归一化后得到
\[
\bigl(\Prob(F_1\mid74),\Prob(F_2\mid74),\Prob(F_3\mid74)\bigr)
=\frac{(100,10,1)}{111}.
\]
较小编号在小批次工厂更常出现, 因而支持第一家来源.

(b) 把“出生前人数”写成文明内从0开始的出生序号. 已见序号小于 $10^{12}$, 所以三个宇宙都容许该观测. 每宇宙大小相同, 先验仍各1/3; 固定序号在文明规模为 $L$ 的宇宙中占个体的比例为 $1/L$. 因而
\[
\Prob(\text{第一宇宙}\mid\text{所见序号})
=\frac{10^{-12}}{10^{-12}+10^{-15}+10^{-18}}
=\frac{1000000}{1001001}.
\]
这个强后验来自给定的均匀抽样及等先验. 对实际文明若这些条件未得到支持, 不能仅凭人类的出生序号套用这一数值.
\end{solution}

\originalsection{作业四: 期望、协方差与效用}
\href{https://ocw.mit.edu/courses/18-600-probability-and-random-variables-fall-2019/resources/mit18_600f19_pset4/}{官方作业四}. 金融与考试故事中的概率均为题设模型, 不是对实际市场或考生成绩的估计.

\paragraph{A 组: 外书题登记.}
Ross 第8版第4章, A.1 为 Problem 23, 含(a)(b). 原件印有题面, 本书仅登记, 不复制外书文字. 原评注关于 Siegel 悖论的链接可作另行阅读.

\begin{exercise}\label{ps4-b}
原题 B. 在 ACT 星球, Jack 每题独立有85\%概率答对, 不受题目对别人难易程度影响. 一次数学 ACT 含60题, 正确51题对应分数30; 要得到34分, 必须至少答对55题. 分数换算每次相同. 大学只要求他至少一次得到34分, 不要求报告低分. 若其能力保持不变且12次考试的所有作答独立, 求至少一次达标的概率, 给代数式和近似百分数.
\end{exercise}
\begin{solution}
一次答对数 $X\sim\Bin(60,0.85)$. 恰对 $j$ 题的概率为 $\binom{60}{j}0.85^j0.15^{60-j}$, 故单次达标率
\[
r=\Prob(X\ge55)=\sum_{j=55}^{60}\binom{60}{j}0.85^j0.15^{60-j}\approx0.0967985087.
\]
12次都不达标的概率为 $(1-r)^{12}$, 所求为
\[
1-\left(1-\sum_{j=55}^{60}\binom{60}{j}0.85^j0.15^{60-j}\right)^{12}
\approx0.7052758311,
\]
即约70.53\%. 能力不变与跨考试独立是将失败概率取12次幂的依据; 若疲劳或共享因素造成关联, 该式就不再直接适用.
\end{solution}

\begin{exercise}\label{ps4-c}
原题 C. 某路口每分钟独立以 $1/500000$ 概率发生事故, 近似一年500000分钟, 因而原年均事故数为1. 新交通灯先验以 $1/3$ 概率有效, 有效时将事故率减半; 以 $2/3$ 概率无效, 保持原率. 安装后的一年没有事故. 用泊松近似求有效的更新概率.
\end{exercise}
\begin{solution}
设 $H$ 为有效, $Z$ 为该年事故数. 有效时近似 $Z\sim\Pois(1/2)$, 无效时近似 $\Pois(1)$. 两者零事故的似然分别为 $\ee^{-1/2}$ 与 $\ee^{-1}$. 因而
\[
\Prob(H\mid Z=0)\approx
\frac{\tfrac13\ee^{-1/2}}{\tfrac13\ee^{-1/2}+\tfrac23\ee^{-1}}
=\frac{\ee^{1/2}}{\ee^{1/2}+2}\approx0.4519.
\]
零事故提升了原来的1/3先验, 但无效时零事故本来就有 $\ee^{-1}$ 的不小概率, 所以并未形成接近确定的证据. 近似针对事故计数的分布, Bayes 更新本身则对这两个泊松模型精确成立.
\end{solution}

\begin{exercise}\label{ps4-d}
原题 D. Larry 发放每笔10000美元的高风险贷款. 25\%的借款人很快只还本金; 50\%失联、破产或死亡而完全不还; 25\%缓慢还款, 连利息与费用共付100000美元, 但 Larry 为催收须向第三方付60000美元.
\begin{enumerate}[label=(\alph*)]
\item 求 Larry 每笔贷款净利润的期望与方差, 净利润须扣除催收开支及最初10000美元支出.
\item 求一个借款人净支付额的期望与方差, 净支付为付出的总额减去最初借得的金额.
\end{enumerate}
这是原件的简化商业故事; 真实借贷法律、道德与行为问题属于其课外评注.
\end{exercise}
\begin{solution}
(a) 净利润 $L$ 在三个情形中依次为 $0,-10000,30000$. 因而
\begin{align*}
\E L&=\tfrac14\cdot0+\tfrac12(-10000)+\tfrac14(30000)=2500,\\
\E L^2&=\tfrac12\cdot10000^2+\tfrac14\cdot30000^2=275000000,\\
\Var(L)&=275000000-2500^2=268750000.
\end{align*}
利润单位是美元, 方差单位是美元平方. 最后一情形利润不是90000, 因为另有60000的催收支出.

(b) 借款人的净支付 $B$ 依次为 $0,-10000,90000$. 未偿还者仍曾得到10000, 所以净支付为负. 计算得
\begin{align*}
\E B&=\tfrac12(-10000)+\tfrac14(90000)=17500,\\
\E B^2&=\tfrac12\cdot10000^2+\tfrac14\cdot90000^2=2075000000,\\
\Var(B)&=2075000000-17500^2=1768750000.
\end{align*}
可用收支平衡再核对: $B-L$ 在三个情形中分别为0、0、60000, 其均值是 $\tfrac14\cdot60000=15000$. 上述计算给 $17500-2500=15000$, 与每笔平均催收费用相符.
\end{solution}

\begin{exercise}\label{ps4-e1}
原题 E.1. 定义 $\Cov(X,Y)=\E(XY)-\E X\E Y$, 并在方差均非零时定义相关系数
$\rho(X,Y)=\Cov(X,Y)/\sqrt{\Var(X)\Var(Y)}$.
核对 $\Cov(X,X)=\Var(X)$, $\Cov(X,Y)=\Cov(Y,X)$, 以及协方差对两个参数分别线性. 例如对实数 $a,b$ 应有 $\Cov(aX+bY,Z)=a\Cov(X,Z)+b\Cov(Y,Z)$.
\end{exercise}
\begin{solution}
假设所有变量有有限二阶矩, 则由 Cauchy--Schwarz, 涉及的乘积可积. 第一式直接为
$\Cov(X,X)=\E X^2-(\E X)^2=\Var(X)$.
交换 $X,Y$ 时乘积及均值乘积都不变, 得到对称性. 再按期望的线性展开,
\begin{align*}
\Cov(aX+bY,Z)
&=\E[(aX+bY)Z]-[a\E X+b\E Y]\E Z\\
&=a[\E XZ-\E X\E Z]+b[\E YZ-\E Y\E Z].
\end{align*}
这就是对第一参数线性; 用对称性得到对第二参数线性. 常数与任意变量的协方差为0, 因此这里的线性不同于“把常数均值相乘”这种误算. 同一中心化表示和 Cauchy--Schwarz 也给 $|\rho|\le1$, 与题面所述范围一致.
\end{solution}

\begin{exercise}\label{ps4-e2}
原题 E.2. 若 $\Cov(X_i,X_j)=ij$, 求 $\Cov(X_1-X_2,X_3-2X_4)$.
\end{exercise}
\begin{solution}
双线性使每一个交叉项都需保留:
\[
\Cov(X_1-X_2,X_3-2X_4)=3-2\cdot4-6+2\cdot8=5.
\]
最后一项为正, 因为两个负系数相乘; 只保留对应位置的两项会漏掉交叉协方差.
\end{solution}

\begin{exercise}\label{ps4-e3}
原题 E.3. 若 $\Cov(X_i,X_j)=ij$, 求 $\Var(X_1+2X_2+3X_3)$.
\end{exercise}
\begin{solution}
把方差视为变量与自身的协方差, 令系数 $a_i=i$, 则
\[
\Var\left(\sum_{i=1}^3iX_i\right)
=\sum_{i=1}^3\sum_{j=1}^3 ij\,\Cov(X_i,X_j)
=\sum_{i,j=1}^3 i^2j^2=(1+4+9)^2=196.
\]
变量之间未给独立, 实际协方差不为0, 所以不能只相加三个单独方差.
\end{solution}

\begin{exercise}\label{ps4-e4}
原题 E.4. 随机变量 $V,X_1,\ldots,X_n$ 满足 $\Var(V)=1$, $\Cov(X_i,V)=b_i$, $\Cov(X_i,X_j)=c_{ij}$, 所有 $b_i,c_{ij}$ 已知. 对固定实系数 $a_i$, 令 $X=\sum_{i=1}^n a_iX_i$. 证明:
\begin{enumerate}[label=(\alph*)]
\item $\Cov(X,V)=\sum_{i=1}^na_ib_i$;
\item $\Var(X)=\sum_{i=1}^n\sum_{j=1}^na_ic_{ij}a_j$;
\item 当 $\Var(X)>0$ 时, $\rho(X,V)=(\sum_i a_ib_i)/\sqrt{\sum_{i,j}a_ic_{ij}a_j}$.
\end{enumerate}
原件提及可进一步用微积分与线性代数优化相关系数, 但未要求完成该优化.
\end{exercise}
\begin{solution}
(a) 固定第二参数 $V$, 逐项用第一参数线性, 得
$\Cov(X,V)=\sum_i a_i\Cov(X_i,V)=\sum_i a_ib_i$.

(b) 两次使用线性, 得
\[
\Var(X)=\Cov\left(\sum_i a_iX_i,\sum_j a_jX_j\right)
=\sum_{i,j}a_i a_j\Cov(X_i,X_j)=\sum_{i,j}a_ic_{ij}a_j.
\]
这也说明右端作为协方差二次型必非负, 但某些非零系数组合仍可能使它为0.

(c) 将(a)(b)和 $\Var(V)=1$ 代入相关系数定义即可. 若该二次型为0, $X$ 几乎处处为常数, 其相关系数没有定义; 不能把 $0/0$ 当作零相关. 原件在相关系数的分式中隐含的非退化条件已在此写明.
\end{solution}

\begin{exercise}\label{ps4-f1}
原题 F.1. 某镇永远每年举行马拉松. 第 $n$ 年男子、女子冠军用时分别为 $M_n,W_n$. 每一序列独立同分布, 分别来自一个连续分布, 分布具体形式无关. 女子第 $n$ 年创造新纪录的概率为 $1/n$, 即 $W_n<W_k$ 对全部 $k<n$ 成立. 为涉及两性联合事件, 以下按原题自然意图取男子序列与女子序列也相互独立; 原件未另用一句话明列这一跨序列条件.
\begin{enumerate}[label=(\alph*)]
\item 求同一年同时创造男子、女子纪录的年数期望.
\item 求素数年份从不同时创造男子、女子纪录的概率.
\item 令 $K\ge2$ 为首次没有女子新纪录的年份, 求 $\E K$. 可用 $\E K=\sum_{k=0}^\infty\Prob(K>k)$.
\item 令 $L$ 为第一年女子冠军的纪录首次被打破的年份. 求 $L$ 为偶数的概率. 原提示为 $\Prob(L\text{偶})=\Prob(L>1)-\Prob(L>2)+\Prob(L>3)-\cdots$.
\item 求满足 $n\ge2$ 且最近两年男子冠军是此前所有冠军中最快两人的年数期望, 即 $\max\{M_n,M_{n-1}\}<M_k$ 对每个 $k<n-1$ 成立.
\end{enumerate}
可使用原件列出的恒等式:
\[
\sum_{n=1}^{\infty}\frac1{n^2}=\frac{\pi^2}{6},\quad
\sum_{n=0}^{\infty}\frac1{n!}=\ee,\quad
\sum_{n=1}^{\infty}\frac{(-1)^{n+1}}n=\ln2,\quad
\sum_{n=2}^{\infty}\binom n2^{-1}=2,
\]
以及 $\sum_{n\ge1}n^s=\prod_{p\text{素数}}(1-p^s)^{-1}$ ($s<-1$). 原件提示应说明纪录事件为什么独立, 不能只引用乘积答案.
\end{exercise}
\begin{solution}
先证明需要的纪录独立性. 连续分布保证任意有限样本几乎必然无并列, 独立同分布又使前 $r$ 年成绩的所有 $r!$ 种相对排序等可能. 记 $U_j$ 为第 $j$ 年成绩在前 $j$ 年中从快到慢的名次. 将第 $j$ 年依名次 $U_j\in\{1,\ldots,j\}$ 插入前 $j-1$ 年的排序, 可以唯一恢复每一种最终排列. 因此名次向量 $(U_1,\ldots,U_r)$ 与 $r!$ 个排列一一对应. 任意合法向量的概率为 $1/r!=\prod_{j=1}^r(1/j)$, 这证明 $U_j$ 相互独立且各自在 $1,\ldots,j$ 上均匀. 创纪录事件正是 $U_j=1$. 男女序列相互独立时, 每年“男女同时创纪录”的指标也跨年独立, 成功率为 $1/j^2$.

(a) 以非负指标求和. 对无穷和使用单调收敛,
\[
\E\sum_{n\ge1}\ind_{\{\text{两纪录}\}}
=\sum_{n\ge1}\frac1{n^2}=\frac{\pi^2}{6}.
\]
第一年自动为两性纪录, 因而 $n=1$ 的1必须包含. 期望有限也说明总次数几乎必然有限.

(b) 先限制到前 $r$ 年素数, 由已证独立性, 没有同时纪录的概率为 $\prod_{p\le r}(1-p^{-2})$. 当 $r\to\infty$, 事件递减, 取极限得到无限乘积. 给定 Euler 乘积恒等式在 $s=-2$ 处成立, 因而
\[
\Prob(\text{所有素数年均无两纪录})
=\prod_p(1-p^{-2})=\left(\sum_{n\ge1}n^{-2}\right)^{-1}=\frac6{\pi^2}.
\]

(c) $K>k$ 表示前 $k$ 年每年都有女子纪录, 即 $W_1>W_2>\cdots>W_k$. 对 $k\ge1$ 这只有 $k!$ 种等可能排序之一, 故概率为 $1/k!$; $k=0$ 时概率也是1. 于是尾和给 $\E K=\sum_{k\ge0}1/k!=\ee$. 尾和来自逐点恒等式 $K=\sum_{k\ge0}\ind_{\{K>k\}}$, 非负性允许交换期望与和.

(d) $L>n$ 表示前 $n$ 年中第一年的成绩最快, 所以对 $n\ge1$, $\Prob(L>n)=1/n$. 由连续性与此尾概率趋零, $L$ 几乎必然有限. 因而
\[
\Prob(L=2j)=\Prob(L>2j-1)-\Prob(L>2j)
=\frac1{2j-1}-\frac1{2j}.
\]
将这些非负项对 $j\ge1$ 求和, 得到成对分组的交错调和级数, 等于 $\ln2$. 配对方式来自概率分解, 不涉及对条件收敛级数的任意重排.

(e) 前 $n$ 年最快两个位置在 $\binom n2$ 个位置对中均匀. 目标只要求这两个位置为 $n-1,n$, 不区分快慢先后, 概率是 $\binom n2^{-1}=2/[n(n-1)]$. 再由非负指标和的单调收敛,
\[
\E(\text{所求年数})=\sum_{n=2}^\infty\frac2{n(n-1)}
=2\sum_{n=2}^\infty\left(\frac1{n-1}-\frac1n\right)=2.
\]
$n=2$ 时条件中没有更早年份, 事件自动成立, 与首项1一致.
\end{solution}

\begin{exercise}\label{ps4-f2}
原题 F.2. 一列无穷长的麻瓜与巫师等候分院. 每次独立以 $1/99$ 概率判为巫师, 然后以相同概率分入4个学院之一; 其他人是麻瓜. 在两次分院之间, 若(a)此前所有人都是巫师, 且(b)4个学院人数相等, 就暂停举办学院间魁地奇赛. 获胜学院的每人得26390金加隆, 学院另得价值1103的奖杯. 开始时4个学院都为0人, 也举办一场空比赛并奖一座奖杯. 求整个过程中将发出的奖励总价值期望. 可用
\[
\frac1\pi=\frac{\sqrt8}{99^2}\sum_{k=0}^{\infty}
\frac{(4k)!(1103+26390k)}{(k!)^4 396^{4k}}.
\]
这里(a)(b)是赛事的两个触发条件, 并不是两道待求子问.
\end{exercise}
\begin{solution}
平衡状态只能出现在处理完 $4k$ 人之后, 且各学院恰有 $k$ 人. 前 $4k$ 人必须全为巫师, 各人的四种巫师学院结果各有概率 $1/(99\cdot4)=1/396$. 恰好每院 $k$ 人的学院序列有多项式系数 $(4k)!/(k!)^4$ 种, 因而赛事发生的概率为
\[
r_k=\frac{(4k)!}{(k!)^4 396^{4k}},\quad k\ge0.
\]
即使同一过程中发生了多场比赛, 各场指标仍可线性相加, 不要求不同 $k$ 的触发事件独立. 每场总奖励为 $1103+26390k$, 因为获胜学院恰有 $k$ 人, 学院获胜机制不影响这一总额. 对非负奖励用单调收敛,
\[
\E(\text{总奖励})=\sum_{k=0}^\infty(1103+26390k)r_k
=\frac{99^2}{\sqrt8\,\pi}.
\]
$k=0$ 项是1103, 正是开始时的奖杯. 另可直接验证期望有限: 平衡事件包含在“前 $4k$ 人全是巫师”内, 所以 $r_k\le99^{-4k}$, 而 $\sum_k(1103+26390k)99^{-4k}<\infty$. 第一次出现麻瓜后不再触发赛事, 也与模型中的概率相符. 所用 Ramanujan 恒等式是原题提供的计算工具, 这里不把它声称为本题中重新证明的结论.
\end{solution}

\begin{exercise}\label{ps4-g1}
原题 G.1. Jill 用二次效用 $U(x)=-(x-x_0)^2$ 评价财富, $x_0$ 为很大的正数. 现有财富 $W_0$, 接受随机净收益 $X$, 其均值为 $\mu$, 方差为 $\sigma^2$, 新财富 $W=W_0+X$. 证明 $\E U(W)$ 仅依赖这两个矩, 并写成 $x_0,W_0,\mu,\sigma^2$ 的函数. 原件说明财富超过 $x_0$ 后效用下降, 暂取 $x_0$ 很大使该情形不太可能.
\end{exercise}
\begin{solution}
令 $c=W_0-x_0$, 展开平方得到
\[
\E U(W)=-\E(c+X)^2=-c^2-2c\mu-(\sigma^2+\mu^2)
=-(W_0+\mu-x_0)^2-\sigma^2.
\]
只用了有限二阶矩, 并未要求 $X$ 服从某种特殊分布. 均值影响平方中的平均位置, 方差单独以负号降低期望效用.
\end{solution}

\begin{exercise}\label{ps4-g2}
原题 G.2. 在上一题模型中, 证明给定收益均值 $\mu$ 后, Jill 偏好尽量小的收益方差 $\sigma^2$. 原件称 $\sigma$ 为风险, 称这种偏好为风险厌恶.
\end{exercise}
\begin{solution}
固定 $W_0,x_0,\mu$ 后, 上题的平方项为常数, 因而两个可选收益的期望效用之差恰为它们方差差的负数. 方差严格更小者效用严格更高, 方差相同者在该模型中无差别. 这是对给定均值的比较; 若收益均值也变, 就须比较完整表达式, 不能只按方差排序.
\end{solution}

\begin{exercise}\label{ps4-g3}
原题 G.3. 若 $X=\sum_{i=1}^na_iX_i$, $a_i$ 为固定实数, $\E X_i=\mu_i$, $\Cov(X_i,X_j)=\sigma_{ij}$, 证明 Jill 的 $\E U(W_0+X)$ 仅依赖这些均值、协方差及系数, 并求表达式. 原提示为先算 $X$ 的均值与方差.
\end{exercise}
\begin{solution}
期望线性与协方差双线性分别给出
\[
\E X=\sum_i a_i\mu_i,\qquad \Var(X)=\sum_{i,j}a_i\sigma_{ij}a_j.
\]
代入 G.1 得
\[
\E U(W_0+X)
=-\left(W_0+\sum_i a_i\mu_i-x_0\right)^2-\sum_{i,j}a_i\sigma_{ij}a_j.
\]
资产间依赖仍通过协方差保留; 我们没有假设独立. 在二次效用的适用范围内, 更高阶联合分布信息不会进入此公式, 这正是原件联系均值--方差投资组合理论的原因.
\end{solution}

\begin{exercise}\label{ps4-g4}
原题 G.4. 设 $X_1,\ldots,X_n$ 独立, 均值相同且方差相同. 在所有 $\sum_{i=1}^na_iX_i$ 中, 系数要求 $a_i\ge0$ 且 $\sum_i a_i=W$, $W$ 为固定常数. 证明等权 $a_i=W/n$ 使方差最小.
\end{exercise}
\begin{solution}
可行性要求 $W\ge0$, 且 $n\ge1$. 设公共方差为 $v\ge0$. 独立和有限二阶矩给不同变量间协方差为0, 所以方差为 $v\sum_i a_i^2$. 由
\[
0\le\sum_{i=1}^n\left(a_i-\frac Wn\right)^2
=\sum_i a_i^2-\frac{W^2}{n},
\]
得方差至少 $vW^2/n$, 等权时达到. 若 $v>0$, 等号当且仅当全部 $a_i=W/n$, 因而最小方案唯一. 若 $v=0$, 所有可行组合的方差都为0, 等权仍最优但不是唯一. $W=0$ 时非负约束已使所有系数为0. 公共均值使各可行方案的收益均值均为同一值, 所以也可用前题的风险厌恶结论比较其效用.
\end{solution}

\originalsection{作业五: 泊松模型}
\href{https://ocw.mit.edu/courses/18-600-probability-and-random-variables-fall-2019/resources/mit18_600f19_pset5/}{官方作业五}. 本节的泊松点过程按不相交时间区间计数独立、区间计数服从相应泊松分布的定义使用. 正式题面在长评注之后仍有 E(f), 该子问也列于下文.

\paragraph{A 组: 外书题登记.}
Ross 第8版第4章, A.1 为 Theoretical Exercise 25, 包括无编号主问及(a)--(c)应用. 原作业公开完整题面, 本书仅保留题号与书目, 不复制外书文字. E(c) 所需泊松细化性质在本书的自设题解中独立推导.

\begin{exercise}\label{ps5-b1}
原题 B.1. $X$ 在 $[0,1]$ 上均匀分布. 对正整数 $n$, 求 $X^n$ 的期望与方差.
\end{exercise}
\begin{solution}
均匀密度在区间内为1, 因而对任意 $r>-1$, $\E X^r=\int_0^1x^r\dd x=1/(r+1)$. 分别取 $r=n,2n$,
\[
\E X^n=\frac1{n+1},\qquad
\Var(X^n)=\frac1{2n+1}-\frac1{(n+1)^2}
=\frac{n^2}{(2n+1)(n+1)^2}.
\]
方差非负且对 $n\ge1$ 为正, 与 $X^n$ 非常数一致. 不能把 $\E X^n$ 写成 $(\E X)^n$, 因为幂不是线性运算.
\end{solution}

\begin{exercise}\label{ps5-b2}
原题 B.2. 若 $X$ 在 $[0,1]$ 上均匀分布, 求 $Y=X^3$ 的概率密度函数.
\end{exercise}
\begin{solution}
立方在此区间单调递增, 所以对 $0\le y\le1$, 有
$F_Y(y)=\Prob(X^3\le y)=\Prob(X\le y^{1/3})=y^{1/3}$.
$y<0$ 时分布函数为0, $y>1$ 时为1. 在 $(0,1)$ 求导得
\[
f_Y(y)=\begin{cases}\tfrac13y^{-2/3},&0<y<1,\\0,&y\notin[0,1].\end{cases}
\]
端点密度值可任意指定, 不影响概率; 0处无原子, 因为 $\Prob(X=0)=0$. 密度在0附近虽无界, 但积分为 $[y^{1/3}]_0^1=1$, 故定义合法.
\end{solution}

\begin{exercise}\label{ps5-c}
原题 C. 规律公交城的班车在 A、B 两站间运行, 无中间站, 每天昼夜在 A 站严格每5分钟发车一次, 如6:00、6:05、6:10等; 一到站即接走所有等车人. 本地居民按每分钟5人的泊松过程到站.
\begin{enumerate}[label=(\alph*)]
\item 你在一天内均匀随机选择的时间到 A 站, 预期还要等多久?
\item 你坐上的车预期载有多少名本地居民?
\end{enumerate}
泊松公交城的班车则按平均每5分钟一辆的泊松点过程昼夜到站, 立即接走所有等车人; 居民到站过程仍为每分钟5人, 且与班车过程独立.
\begin{enumerate}[label=(\alph*),start=3]
\item 你仍在一天内均匀随机选择的时间到站, 预期等多久?
\item 你坐上的车预期有多少名本地居民?
\item 下列两句话是真是假? 若均真, 用文字解释看似矛盾之处: (i) 对随机到访的你而言, 泊松城公交看起来平均等候时间为规律城的两倍, 且车上为两倍拥挤; (ii) 两城长期都是平均每5分钟一辆车、每分钟5人到站, 因而长期平均每车都是25人, 一样拥挤.
\end{enumerate}
这里的居民数不包括作为额外访客的你. 原件还讨论公交成对到达造成更严重的等待与拥挤, 但该评注不另设子问.
\end{exercise}
\begin{solution}
(a) 一天含整数个5分钟周期. 均匀到达时, 在当前周期的位置均匀, 剩余等待 $R$ 也在 $[0,5]$ 上均匀, 所以 $\E R=2.5$ 分钟. 恰好与公交同时到达的事件概率为0, 不影响结果.

(b) 你所乘车将收取前一班车到这班车的整个5分钟内的居民, 包括你到站前已经等待的和之后到站的. 居民到站过程独立于访客时间, 该固定区间计数服从 $\Pois(25)$, 所以期望为25人. 只算自己等待期间新到的人, 会漏掉已有乘客.

(c) 为体现昼夜稳定运行, 取班车过程为实轴上的平稳泊松过程, 率 $\lambda=1/5$ 每分钟. 固定一个与该过程独立的访客时刻. 未来 $r$ 分钟无班车的概率为 $\ee^{-r/5}$, 故前向等待 $R\sim\Exp(1/5)$, $\E R=5$ 分钟. 均匀随机选当天时刻不改变这一条件分布.

(d) 从访客时刻向后看, 上一辆车到现在的时间 $A$ 也有相同指数分布. 过去与未来的区间计数独立, 所以 $A,R$ 独立, $\E(A+R)=10$ 分钟. 你所乘车的收客区间长为 $A+R$. 条件于此区间及公交过程, 居民计数的期望为 $5(A+R)$, 再取期望得50人.

(e) 两句在上述解释下均真. 对随机时刻观察者, 一个长度10分钟的公交间隔被遇到的机会为长度5分钟间隔的两倍; 这种按时间抽样偏重长间隔. 因而访客所处间隔平均长10分钟, 其结束处公交平均收50人. 逐辆公交抽样则每辆只计一次, 普通相邻车间隔平均长5分钟, 条件乘客均值是其长度的5倍, 平均为25人. 两个平均对应不同抽样方式, 并未对同一随机对象给相反答案. 长期而言, 独立指数间隔及各间隔的泊松居民计数构成同分布的周期, 用大数定律得到每车平均25人; 与“随机访客坐到哪辆车”的长度偏差完全相容.
\end{solution}

\begin{exercise}\label{ps5-d}
原题 D. 每一天独立于其余天, Jill 以 $1/2500$ 概率听到某一事实, 如 Henry Mancini 创作了《粉红豹主题曲》. 听到累计3次后才能进入长期记忆. Alice 每天以 $1/1000$ 概率听到同一事实, 听到2次后记住. 使用泊松近似回答:
\begin{enumerate}[label=(\alph*)]
\item Jill 在10000天大时已记住该事实的概率;
\item Alice 在10000天大时已记住该事实的概率;
\item 若有10000个具有同样接触、记忆概率的事实, 且对每个给定事实, Jill 与 Alice 的听到次数相互独立, 两人均为10000天大时, 预期有多少事实 Jill 知道而 Alice 不知道?
\end{enumerate}
\end{exercise}
\begin{solution}
(a) Jill 的二项接触计数均值为 $10000/2500=4$, 近似 $J\sim\Pois(4)$. 所求是听到至少3次, 所以
\[
\Prob(J\ge3)\approx1-\ee^{-4}\left(1+4+\frac{4^2}{2}\right)
=1-13\ee^{-4}\approx0.7618966944.
\]

(b) Alice 的计数均值为10, 近似 $A\sim\Pois(10)$. 她记住的概率为
$\Prob(A\ge2)\approx1-\ee^{-10}(1+10)=1-11\ee^{-10}\approx0.9995006008$.

(c) 对某一个事实, Jill 知道且 Alice 不知道的概率, 按题给的二人独立性, 为
$(1-13\ee^{-4})11\ee^{-10}$.
把10000个事实的指标相加, 期望为
\[
10000(1-13\ee^{-4})11\ee^{-10}\approx3.8049062.
\]
不同事实之间是否独立不影响这一期望, 因为使用的是期望线性; 同一事实在两人之间的独立性则用于乘积, 确实不可省略. 这些数值全部是按题意的泊松近似, 不是精确二项尾概率.
\end{solution}

\begin{exercise}\label{ps5-e}
原题 E. 在 $[0,\infty)$ 上取率1的泊松点过程, 到达时刻记为 $X_1<X_2<\cdots$. $X_1,X_2-X_1,\ldots$ 是独立的率1指数变量. 令 $Y_i=\ln X_i$. 在“到达频率加速星球”(AFP), 公交到达时刻为 $Y_i$, 时间从 $-\infty$ 延续到 $\infty$, 因而可能有负的到达时刻.
\begin{enumerate}[label=(\alph*)]
\item 对任意 $a<b$, 证明时间区间 $[a,b]$ 的公交数为参数 $\int_a^b\ee^x\dd x$ 的泊松变量.
\item 各用一句话说明以下事实: $(-\infty,0]$ 的公交数为参数1的泊松变量; 零时以后几乎必然有无穷辆公交; 若 $\alpha=\ln2\approx0.7$, 则 $[k\alpha,(k+1)\alpha]$ 的预期公交数为 $[(k-1)\alpha,k\alpha]$ 的两倍; 首辆公交时刻的中位数为 $\ln\ln2\approx-0.37$.
\end{enumerate}
AFP 一只不能移动的树懒在时刻 $-7$ 生于车站, 直到出生后第一辆会杀死它的公交到达. 标准时间单位为十二年, 故 $\alpha$ 单位约8.32年. 因出生前有公交的概率很小, 原件用首辆时刻中位数估计树懒寿命中位数约 $12(7-0.37)\approx80$ 年.
\begin{enumerate}[label=(\alph*),start=3]
\item 每辆公交独立掷公平硬币决定为粗轮胎或细轮胎. 女性树懒只被粗轮胎公交杀死, 男性被任意公交杀死. 说明女性的预期寿命约比男性长8.32年, 即 $\alpha$ 个时间单位. 原提示用泊松细化及出生时刻平移, 并注意 $-7-\alpha$ 到 $-7$ 之间有公交的概率很小.
\item 令 $p_k$ 为已存活 $k-1$ 年的树懒在第 $k$ 年死亡的条件概率. 说明当 $p_k$ 很小, 比如小于0.1时, 它近似等于该年预期公交数, 因而前约80年随年龄大致指数增长.
\item 查阅 \href{https://www.ssa.gov/oact/STATS/table4c6.html}{SSA 生命表} 与 \href{https://en.wikipedia.org/wiki/Gompertz%E2%80%93Makeham_law_of_mortality}{Gompertz--Makeham 死亡律介绍}, 花约5分钟观察后写一两句话, 特别比较人类年度死亡概率与 AFP 树懒模型的接近程度. 原题指出三个明显差别: 人类婴幼儿的死亡率较高; 十几岁末至二十多岁、尤其男性, 出现高于简单 Gompertz 趋势的凸起; 其他年龄段的人类男女死亡率差距较大但小于本树懒模型. 可另读原件所列 Gravity and Levity 博客的说明.
\item 在另一个严格 Gompertz 世界中, 死亡强度每8年加倍, 不同人的寿命相互独立. 四兄弟姐妹当前年龄为20、28、28、36岁. 求最年长者最先死亡的概率.
\end{enumerate}
(f) 使用原件后段评注中的严格8年倍增世界, 并非把前面8.32年的近似偷偷替换掉.
\end{exercise}
\begin{solution}
(a) 单调变换 $x=\ee^y$ 将 AFP 区间 $[a,b]$ 对应为原过程区间 $[\ee^a,\ee^b]$. 原过程区间计数为 $\Pois(\ee^b-\ee^a)$, 且端点命中概率为0. 因 $\ee^b-\ee^a=\int_a^b\ee^y\dd y$, 所证成立. 不相交 AFP 区间对应不相交原区间, 因而仍有独立增量; 强度随时间为 $\lambda(y)=\ee^y$.

(b) 第一事实: $Y\le0$ 等价于 $X\le1$, 对应原过程长度1的区间. 第二事实: 原过程在 $[1,\infty)$ 几乎必然无限到达; 例如各单位区间非空事件独立且概率恒为 $1-\ee^{-1}$, 对任意起点其后全空的概率为0, 对整数起点作可数并即可. 第三事实: $[k\alpha,(k+1)\alpha]$ 的均值为 $\ee^{k\alpha}(\ee^\alpha-1)=2^k$, 上一相邻区间均值为 $2^{k-1}$. 第四事实: $\Prob(Y_1\le y)=\Prob(X_1\le\ee^y)=1-\ee^{-\ee^y}$, 令其为1/2得 $y=\ln\ln2$. 这四句的概率依据都来自时间变换, 不是把等间隔公交当成指数过程.

(c) 先独立证明所用细化性质. 对一个原计数 $N\sim\Pois(\Lambda)$, 每辆以概率 $p$ 标为致命, 则致命数 $K$ 条件于 $N=n$ 服从 $\Bin(n,p)$. 对 $k\ge0$,
\begin{align*}
\Prob(K=k)
&=\sum_{n=k}^\infty\ee^{-\Lambda}\frac{\Lambda^n}{n!}\binom nkp^k(1-p)^{n-k}\\
&=\ee^{-\Lambda}\frac{(\Lambda p)^k}{k!}
\sum_{m=0}^\infty\frac{[\Lambda(1-p)]^m}{m!}
=\ee^{-p\Lambda}\frac{(p\Lambda)^k}{k!}.
\end{align*}
区间计数与各车标记独立, 所以细化后不相交区间仍独立. 因此女性致命公交强度为 $\ee^y/2=\ee^{y-\alpha}$.

设出生时刻 $a=-7$, 以十二年为单位的剩余寿命为 $T$. 男性存活到出生后 $t\ge0$ 的概率, 是区间 $(a,a+t]$ 无公交的概率:
\[
S_{m,a}(t)=\exp\{-\ee^a(\ee^t-1)\}.
\]
女性相应为 $S_{f,a}(t)=\exp\{-\tfrac12\ee^a(\ee^t-1)\}=S_{m,a-\alpha}(t)$, 所以女性在 $a$ 出生的寿命长度分布精确等于男性在 $a-\alpha$ 出生的长度分布. 还需解释为何这使平均长度近似增加 $\alpha$: 向前平移出生时刻通常只是多等 $\alpha$, 但若这新增时间内已有公交, 早出生男性就会先死, 所以不是严格加常数.

精确地, 令 $c=\ee^a$. 尾积分给出男性期望寿命
\[
g(c)=\int_0^\infty\ee^{-c(\ee^t-1)}\dd t
=\ee^c\int_c^\infty\frac{\ee^{-u}}u\dd u.
\]
女性期望为 $g(c/2)$. 为控制小 $c$ 时的误差, 写
\[
\int_c^\infty\frac{\ee^{-u}}u\dd u
=-\ln c+C+r(c),\qquad
r(c)=\int_0^c\frac{1-\ee^{-u}}u\dd u,
\]
其中常数 $C=\int_0^1(\ee^{-u}-1)\dd u/u+\int_1^\infty\ee^{-u}\dd u/u$ 有限, 且 $0\le r(c)\le c$. 因此 $g(c)=-\ln c+C+O(c|\ln c|)$, 得 $g(c/2)-g(c)=\ln2+O(c|\ln c|)$. 对 $c=\ee^{-7}$ 很小, 差值接近 $\alpha$.

按“首辆公交时刻”直接平移的误差也可从新增区间看出: $(a-\alpha,a]$ 的期望公交数为 $c/2$, 至少一辆的概率为 $1-\ee^{-c/2}<c/2\approx0.000456$. 用上述精确尾积分数值核对, 男、女平均寿命约77.15、85.44年, 差约8.28年; 忽略出生边界的小修正即为原题的 $12\ln2\approx8.32$ 年. 这些是树懒模型的数学均值, 不应与前面约80年的中位数混淆.

(d) 为对应题面, 先算被所有公交杀死的男性. 第 $k$ 年对应 AFP 时间区间
$[a+(k-1)/12,a+k/12]$. 在已存活至其起点的条件下, 未来增量仍独立, 所以
\[
\Lambda_k=\ee^{-7}\ee^{(k-1)/12}(\ee^{1/12}-1),\qquad
p_k=1-\ee^{-\Lambda_k}.
\]
若女性, 将 $\Lambda_k$ 减半即可. 对 $\Lambda\ge0$, Taylor 余项或积分估计给
$0\le\Lambda-(1-\ee^{-\Lambda})\le\Lambda^2/2$.
因此小强度时 $p_k\approx\Lambda_k$, 并且 $\Lambda_{k+1}/\Lambda_k=\ee^{1/12}$, 随年数指数增长. 当 $p_k<0.1$, $\Lambda_k<-\ln0.9\approx0.10536$, 相对误差约在5\%量级以内. 例如男性第80年的模型死亡概率约0.05568, 仍在所述小概率区间. 很高龄时 $p_k$ 接近1, 不会永远作为概率指数增长; 真正严格指数增长的是强度.

(e) 这是资料观察题, 没有唯一数值答案. 一个可复核的比较方法是从所选年份生命表取年龄 $k-1$ 的一年死亡概率 $q_{k-1}$, 转为累计强度 $h_k=-\ln(1-q_{k-1})$, 再比较 $\ln h_k$ 是否随年龄近似线性及其男女差. AFP 模型有 $h_k=\Lambda_k$, 故图上严格是一条直线, 女性的纵坐标比男性低 $\ln2$. 小 $q$ 时也可直接比较 $\ln q$, 但接近1时这种近似会失效.

依据原题给出的观察提示, 一两句示例回答可为: “简单指数模型较适合描述成人中间年龄段的上升趋势, 但不能同时解释婴幼儿的较高死亡率和青年期的凸起; 男女曲线的间距也不能直接取树懒模型中固定的因子二.” 这是一份依据原题提示的定性回答示例, 不是新拟合出来的实证结论. 生命表网页会更新, 作实际观察时应注明表年、查阅日期和所比较年龄段; 不把现时网页数据改写成2019年题面中的历史数字. 此处没有声称已重建原作业当时的生命表快照.

本次于2026年10月8日打开上述SSA入口时, 所显示的是2023时期生命表, 用于2026年 Trustees Report; 以下两句只完成当前页面的观察, 不替代2019年题面或声称复原其快照. 该表0、20、60岁的男性一年死亡概率依次为0.006015、0.001235、0.011337, 女性依次为0.005125、0.000441、0.006923, 可见出生首年高于20岁, 而成年后到60岁又明显升高. 20岁与60岁的男女死亡概率比分别约2.80与1.64, 并非AFP模型在小死亡概率时近似二倍的关系(其累计强度比精确为二); 青年期不能据此简单断言所有年龄段的性别差均小于二倍.

(f) 改用严格8年倍增世界. 以20岁者当前强度为共同基准, 四人的强度权重分别为 $1,2,2,4$. 未来共同年龄增长使它们都乘同一函数 $h(t)$, 所以最年长者强度为 $4h(t)$, 总强度为 $9h(t)$. 不同人寿命独立, 尚无人死亡至 $t$ 的概率为
$\exp\{-9\int_0^t h(s)\dd s\}$.
故最年长者最先死的概率为
\[
\int_0^\infty4h(t)\exp\left\{-9\int_0^th(s)\dd s\right\}\dd t
=\frac49\int_0^\infty9h(t)\ee^{-9\int_0^t h(s)\dd s}\dd t=\frac49.
\]
最后积分为1, 因 Gompertz 强度的累计积分趋于无穷, 四人中几乎必然有首个死亡, 且连续分布无并列. 权重比决定谁最先死, 不意味着年龄最大者一定最先死.
\end{solution}
